% ================================
% Packages
% ================================
\usepackage[francais]{babel}
\usepackage[letterpaper,top=2cm,bottom=2cm,left=3cm,right=3cm,marginparwidth=1.75cm]{geometry}
\usepackage{amsmath,amsthm,amssymb,amsfonts,stmaryrd,bbold,graphicx,tcolorbox,tikz,mathabx,mathtools,fourier-orns,enumitem} % Required for inserting images
\usepackage[T1]{fontenc}
\usepackage[dvipsnames]{xcolor}
\usepackage[hidelinks]{hyperref}
\usepackage[nottoc]{tocbibind}
\usepackage{pgfplots}
\pgfplotsset{compat=1.18}


% ================================
% Style "poly"
% ================================

\renewcommand{\thesection}{\Roman{section}}

\frenchbsetup{StandardLists=true}
\tcbuselibrary{%
  theorems,%
  breakable,%
  skins
}
\usetikzlibrary{automata, positioning,fadings,shapes,arrows.meta,calc,fit,backgrounds,shapes.multipart,positioning,patterns,decorations.pathmorphing} 

\tcbset{set1/.style={%
  colback=white,%
  separator sign colon,%
  description delimiters none,%
  enhanced jigsaw,%
  breakable,%
  sharp corners=all,
  coltitle=black,
  frame hidden,
  borderline west={1pt}{0pt}{black},
  overlay unbroken and last={%
        \draw[line width=1pt, black] (frame.south west)   -- ++(0:1cm);},
  sharp corners=southwest,
  before skip=10pt, after skip=10pt,
  left=6pt, right=6pt, top=6pt, bottom=6pt,
attach boxed title to top left={xshift=1.5mm},
boxed title style={skin=enhanced jigsaw,
colback=white,sharp corners,frame hidden,  borderline south={1pt}{0pt}{black},left=-3pt,right=-3pt,bottom=-1pt
}}
}







\newtcbtheorem[number within=section]{df}{Définition}{set1}{df}
\newtcbtheorem[number within=section]{re}{Remarque}{set1}{re}
\newtcbtheorem[number within=section]{prop}{Proposition}{set1}{prop}
\newtcbtheorem[use counter from=df]{dfprop}{Définition-Proposition}{set1}{dfprop}
\newtcbtheorem[number within=section]{tm}{Théorème}{set1}{tm}
\newtcbtheorem[number within=section]{lm}{Lemme}{set1}{lm}
\newtcbtheorem[number within=section]{ex}{Exemple}{set1}{ex}
\newtcbtheorem[number within=section]{exo}{Exercice}{set1}{exo}
\newtcbtheorem[number within=section]{cor}{Corollaire}{set1}{cor}
\newtcbtheorem[number within=section]{ra}{Rappel}{set1}{ra}
\newtcbtheorem[number within=section]{nt}{Notation}{set1}{nt}
\newtcbtheorem{qs}{Question}{set1}{qs}

% ================================
% Nouveaux opérateurs et définitions
% ================================

\definecolor{treeblue}{RGB}{50, 90, 180}
\definecolor{labelgreen}{RGB}{0, 140, 130}
\definecolor{gridblue}{RGB}{70, 100, 200}
\definecolor{signpink}{RGB}{220, 60, 150}
\definecolor{myblue}{RGB}{50, 90, 180}
\definecolor{mypink}{RGB}{180, 50, 130}
\DeclareMathOperator{\inde}{\perp \!\!\! \perp}

% ================================
% Raccourcis
% ================================

\newcommand{\ssi}{\text{si et seulement si} }
\newcommand{\tq}{\text{tel que} }
\newcommand{\apcr}{\text{à partir d'un certain rang} }
\newcommand{\ev}{\text{espace vectoriel} }
\newcommand{\evs}{\text{espaces vectoriels} }
\newcommand{\sev}{\text{sous espace vectoriel} }
\newcommand{\sevs}{\text{sous espaces vectoriels} }
\newcommand{\pgcd}{\text{pgcd}}
\newcommand{\ppcm}{\text{ppcm}}
\newcommand{\evn}{\text{espace vectoriel normé} }
\newcommand{\gr}{\text{groupe} }
\newcommand{\grs}{\text{groupes} }
\newcommand{\sg}{\text{sous-groupe} }
\newcommand{\sgs}{\text{sous-groupes} }
\newcommand{\sgd}{\text{sous-groupe distingué} }
\newcommand{\sgds}{\text{sous-groupes distingués} }
\newcommand{\aut}{\text{Aut}}
\def\Xint#1{\mathchoice
{\XXint\displaystyle\textstyle{#1}}%
{\XXint\textstyle\scriptstyle{#1}}%
{\XXint\scriptstyle\scriptscriptstyle{#1}}%
{\XXint\scriptscriptstyle\scriptscriptstyle{#1}}%
\!\int}
\def\XXint#1#2#3{{\setbox0=\hbox{$#1{#2#3}{\int}$}
\vcenter{\hbox{$#2#3$}}\kern-.5\wd0}}
\def\dashint{\Xint-}

\newcommand{\toninfty}{\underset{n\to +\infty}{\longrightarrow}}
\newcommand{\tops}{\overset{\text{p.s.}}{\longrightarrow}}
\newcommand{\toproba}{\overset{\text{proba}}{\longrightarrow}}
\newcommand{\support}{\text{supp}}
\newcommand{\Lloc}{L^1_{\text{loc}}}
\newcommand{\1}{\mathbb{1}}
\newcommand{\N}{\mathbb{N}}
\newcommand{\Z}{\mathbb{Z}}
\newcommand{\R}{\mathbb{R}}
\newcommand{\Rb}{\overline{\R}}
\newcommand{\Pa}{\mathbb{P}}
\newcommand{\Q}{\mathbb{Q}}
\newcommand{\C}{\mathbb{C}}
\newcommand{\K}{\mathbb{K}}
\newcommand{\Fa}{\mathbb{F}}
\newcommand{\La}{\mathbb{L}}
\newcommand{\Ea}{\mathbb{E}}
\newcommand{\Gr}{\mathcal{G}}
\newcommand{\Kr}{\mathcal{K}}
\newcommand{\Ar}{\mathcal{A}}
\newcommand{\Br}{\mathcal{B}}
\newcommand{\Cr}{\mathcal{C}}
\newcommand{\Mr}{\mathcal{M}}
\newcommand{\Hr}{\mathcal{H}}
\newcommand{\Nr}{\mathcal{N}}
\newcommand{\Fr}{\mathcal{F}}
\newcommand{\Ur}{\mathcal{U}}
\newcommand{\Yr}{\mathcal{Y}}
\newcommand{\Zr}{\mathcal{Z}}
\newcommand{\Lr}{\mathcal{L}}
\newcommand{\Or}{\mathcal{O}}
\newcommand{\Rr}{\mathcal{R}}
\newcommand{\Dr}{\mathcal{D}}
\newcommand{\Er}{\mathcal{E}}
\newcommand{\Par}{\mathcal{P}}
\newcommand{\Tr}{\mathcal{T}}
\newcommand{\Vr}{\mathcal{V}}
\newcommand{\Sr}{\mathcal{S}}
\newcommand{\Xr}{\mathcal{X}}
\newcommand{\BR}{\Br(\R)}
\newcommand{\BRb}{\Br(\Rb)}
\newcommand{\RBR}{(\R,\BR)}
\newcommand{\RBRb}{(\Rb,\BRb)}
\newcommand{\nN}{{n\in\N}}
\newcommand{\wt}{\backslash}
\newcommand{\lims}{\overline{\lim}}
\newcommand{\limi}{\underline{\lim}}
\newcommand{\EA}{(E,\Ar)}
\newcommand{\EAmu}{(E,\Ar,\mu)}
\newcommand{\FB}{(F,\Br)}
\newcommand{\FBnu}{(F,\Br,\nu)}
\newcommand{\LEA}{L^1(E,\Ar,\mu)}
\newcommand{\LFB}{L^1(F,\Br,\nu)}
\newcommand{\OmegaFP}{(\Omega,\Fr,\Pa)}
\newcommand{\LCEA}{L_\C^1(E,\Ar,\mu)}
\newcommand{\pp}[1]{#1\text{-presque partout} }
\newcommand{\va}{\text{variable aléatoire} }
\newcommand{\vas}{\text{variables aléatoires} }
\newcommand{\LRBR}{L^1(\R,\BR,\lambda)}
\newcommand{\cl}{{\text{cl}}}
\newcommand{\Sn}{\Sr_n}
\newcommand{\An}{\Ar_n}
\newcommand{\id}{\text{id} }
\newcommand{\Id}{\text{Id} }
\newcommand{\Var}{\text{Var} }
\newcommand{\Cov}{\text{Cov} }
\newcommand{\Hom}{\text{Hom} }
\newcommand{\car}{\text{car} }
\newcommand{\Aut}{\text{Aut} }
\newcommand{\Gal}{\text{Gal} }
\newcommand{\iid}{\text{i.i.d.} }
\newcommand{\sinc}{\text{sinc} }

% ================================
% Raccourcis à paramètres
% ================================

\newcommand{\produit}[2]{\prod_{#1}^{#2}}
\newcommand{\vect}[1]{\overrightarrow{#1}}
\newcommand{\bnorme}[1]{ \left\lVert #1 \right\rVert }
\newcommand{\norme}[1]{ \lVert #1 \rVert }
\newcommand{\acc}[1]{ \left\{ #1 \right\} }
\newcommand{\cho}[1]{ \left\{\begin{array}{ll} #1 \end{array} \right. }
\newcommand{\fun}[2]{\begin{array}{ccccc}
#1 & \to & #2
\end{array}}
\newcommand{\funinj}[2]{\begin{array}{ccccc}
#1 & \hookrightarrow & #2
\end{array}}
\newcommand{\funsurj}[2]{\begin{array}{ccccc}
#1 & \twoheadrightarrow & #2
\end{array}}
\newcommand{\mfun}[2]{\begin{array}{ccccc}
#1 & \mapsto & #2
\end{array}}
\newcommand{\funto}[3]{\begin{array}{ccccc}
#1 : & #2 & \to & #3
\end{array}}
\newcommand{\funtoinj}[3]{\begin{array}{ccccc}
#1 : & #2 & \hookrightarrow & #3
\end{array}}
\newcommand{\funtosurj}[3]{\begin{array}{ccccc}
#1 : & #2 & \twoheadrightarrow & #3
\end{array}}
\newcommand{\fundf}[4]{\begin{array}{ccccc}
& #1 & \to & #2 \\
& #3 & \mapsto & #4  
\end{array}}
\newcommand{\fundfinj}[4]{\begin{array}{ccccc}
& #1 & \hookrightarrow & #2 \\
& #3 & \mapsto & #4  
\end{array}}
\newcommand{\fundfsurj}[4]{\begin{array}{ccccc}
& #1 & \twoheadrightarrow & #2 \\
& #3 & \mapsto & #4  
\end{array}}
\newcommand{\fundfto}[5]{\begin{array}{ccccc}
#1 : & & #2 & \to & #3 \\
& & #4 & \mapsto & #5  
\end{array}}
\newcommand{\fundftoinj}[5]{\begin{array}{ccccc}
#1 : & & #2 & \hookrightarrow & #3 \\
& & #4 & \mapsto & #5  
\end{array}}
\newcommand{\fundftosurj}[5]{\begin{array}{ccccc}
#1 : & & #2 & \twoheadrightarrow & #3 \\
& & #4 & \mapsto & #5  
\end{array}}
\newcommand{\sca}[1]{\langle #1 \rangle}
\newcommand{\ext}[2]{\begin{matrix}
  #1\\
  |\\
  #2
\end{matrix}}